\section*{Bab \ref{ch:g12:space} --- Vektor, Garis dan Bidang di Ruang}

\begin{solution}{exo:g12:space:1}
$\vect{AB}(2, 1, -1)$ dan $\vect{AC}(1, -1, 1)$ tidak sebanding
($\frac21 \neq \frac{1}{-1}$), jadi titiknya tidak \omterm{def:g12:space:collinear}{segaris}. Garis $(AB)$:
\[
(x, y, z) = (1 + 2t,\; t,\; 2 - t), \qquad t \in \R .
\]
\end{solution}

\begin{solution}{exo:g12:space:2}
$2(x - 1) - (y - 1) + 3(z - 0) = 0$, \emph{yaitu}
\[
2x - y + 3z - 1 = 0 .
\]
Untuk $B(0,2,1)$: $0 - 2 + 3 - 1 = 0$, jadi ya, $B$ terletak pada bidangnya.
\end{solution}

\begin{solution}{exo:g12:space:3}
$\vect{AB}(1,1,0)$, $\vect{AC}(1,0,1)$:
\[
\cos\widehat{A}
= \frac{\vect{AB}\cdot\vect{AC}}{\norm{\vect{AB}}\,\norm{\vect{AC}}}
= \frac{1}{\sqrt2\,\sqrt2} = \frac12,
\]
jadi $\widehat A = 60^\circ$.
\end{solution}

\begin{solution}{exo:g12:space:4}
Titik pada $\mathcal D$: $(1 + t,\, 2 - t,\, 2t)$. Dengan mensubstitusikannya ke
\omterm{def:g10:algebra:equation}{persamaan} $\mathcal P$:
\[
2(1+t) + (2 - t) - 2t + 1 = 5 - t = 0 \iff t = 5 .
\]
Titik potongnya tunggal: $(6, -3, 10)$.
\end{solution}

\begin{solution}{exo:g12:space:5}
\omterm{def:g12:space:normal}{Vektor normalnya} $(1, -1, 2)$ dan $(2, 1, 1)$ tidak \omterm{def:g12:space:collinear}{segaris}, sehingga kedua
bidangnya berpotongan pada sebuah garis. Dengan menjumlahkan kedua \omterm{def:g10:algebra:equation}{persamaannya}:
$3x + 3z = 5$, jadi $x = \frac53 - z$; lalu \omterm{def:g10:algebra:equation}{persamaan} pertamanya memberi
$y = x + 2z - 1 = \frac23 + z$. Dengan $z = t$:
\[
(x, y, z) = \left(\frac53 - t,\; \frac23 + t,\; t\right), \qquad t \in \R,
\]
yaitu garis dengan \omterm{def:g11:vect:direction}{vektor arah} $(-1, 1, 1)$. (Kedua \omterm{def:g10:algebra:equation}{persamaannya} lolos
pemeriksaan.)
\end{solution}

\begin{solution}{exo:g12:space:6}
Arahnya $\vec u_1(1, 2, -1)$ dan $\vec u_2(1, 1, 2)$ tidak \omterm{def:g12:space:collinear}{segaris},
jadi kedua garisnya tidak sejajar. Agar berpotongan diperlukan
\[
1 + t = 2 + s, \qquad 2t = 1 + s, \qquad 3 - t = 1 + 2s .
\]
Dua yang pertama memberi $t = 1 + s$ dan $2(1+s) = 1 + s$, jadi $s = -1$, $t = 0$;
tetapi yang ketiga lalu terbaca $3 = -1$, dan itu mustahil. Tanpa titik persekutuan
dan tidak sejajar: jadi kedua garisnya bersilangan.
\end{solution}

\begin{solution}{exo:g12:space:7}
\emph{1.} $\operatorname{dist} = \dfrac{\abs{2\cdot3 - 1 + 2\cdot1 - 3}}
{\sqrt{4 + 1 + 4}} = \dfrac{4}{3}$.

\emph{2.} $H = M_0 + t\,\vec n$ dengan $\vec n(2, -1, 2)$, dan $t$ dipilih agar
$H \in \mathcal P$:
\[
2(3 + 2t) - (1 - t) + 2(1 + 2t) - 3 = 4 + 9t = 0 \iff t = -\frac49 .
\]
Jadi $H\left(\frac{19}{9}, \frac{13}{9}, \frac19\right)$, dan
\[
M_0H = \norm{t\,\vec n} = \frac49 \times 3 = \frac43,
\]
yang cocok dengan rumus jaraknya.
\end{solution}

\begin{solution}{exo:g12:space:8}
\omterm{def:g10:coordgeom:system}{Koordinatnya}: $G(1,1,1)$, dan bidang $(BDE)$ lewat $B(1,0,0)$,
$D(0,1,0)$, $E(0,0,1)$.

\emph{1.} Bidang $(BDE)$ berpersamaan $x + y + z = 1$ (dipenuhi ketiga
titiknya, dan ketiganya tidak \omterm{def:g12:space:collinear}{segaris}), sehingga $\vec n(1,1,1)$ normal
terhadapnya. Karena $\vect{AG} = (1,1,1) = \vec n$, maka diagonal $(AG)$ \omterm{def:g11:scal:orthogonal}{ortogonal}
terhadap bidang $(BDE)$.

\emph{2.} Jarak dari $A(0,0,0)$:
$\frac{\abs{0 + 0 + 0 - 1}}{\sqrt3} = \frac{1}{\sqrt3} = \frac{\sqrt3}{3}$.
Garis $(AG)$ adalah $(t, t, t)$; garis itu menemui bidangnya ketika $3t = 1$, yaitu di
titik $\left(\frac13, \frac13, \frac13\right)$, yang merupakan titik berat
segitiga $BDE$.
\end{solution}

\begin{solution}{exo:g12:space:9}
\emph{1.} $\vect{AB}(1, 0, -1)$, $\vect{AC}(-1, 1, 0)$,
$\vect{AD}(1, 1, 1)$. Jika $\vect{AD} = a\vect{AB} + b\vect{AC}$, maka
\omterm{def:g10:coordgeom:system}{koordinatnya} memberi $a - b = 1$, $b = 1$, $-a = 1$: dua yang pertama memaksa
$a = 2$, dan itu bertentangan dengan $a = -1$. Jadi tidak \omterm{def:g12:space:collinear}{sebidang}.

\emph{2.} $\vect{BC}(-2, 1, 1)$, $\vect{BD}(0, 1, 2)$. Dengan menyelesaikan
$\vec n \cdot \vect{BC} = \vec n \cdot \vect{BD} = 0$: dari
$b + 2c = 0$ diperoleh $b = -2c$; lalu $-2a - 2c + c = 0$ memberi $c = -2a$. Dengan
$a = 1$: $\vec n(1, 4, -2)$. Bidang lewat $B(2,1,0)$:
\[
x + 4y - 2z - 6 = 0 .
\]

\emph{3.} Jarak dari $A(1,1,1)$:
$\dfrac{\abs{1 + 4 - 2 - 6}}{\sqrt{21}} = \dfrac{3}{\sqrt{21}}$.
Luas $BCD$: $\norm{\vect{BC}}^2 = 6$, $\norm{\vect{BD}}^2 = 5$,
$\vect{BC}\cdot\vect{BD} = 3$, sehingga
\[
\text{Luas} = \frac12\sqrt{6 \times 5 - 9} = \frac{\sqrt{21}}{2}.
\]

\emph{4.}
\[
V = \frac13 \times \frac{\sqrt{21}}{2} \times \frac{3}{\sqrt{21}}
= \frac12 .
\]
\end{solution}

\begin{solution}{pb:g12:space:1}
\textbf{1.} Bidangnya: $x + y + z = 1$. Garisnya: $(t, t, t)$.
Titik potongnya: $3t = 1$, jadi
$\left(\frac13, \frac13, \frac13\right)$.

\textbf{2.} $\dfrac{\abs{0 + 0 + 0 - 1}}{\sqrt3} =
\dfrac{1}{\sqrt3} = \dfrac{\sqrt3}{3}$.

\textbf{3.} $1 - 1 + 0 = 0$: jadi \omterm{def:g11:scal:orthogonal}{ortogonal}.

\textbf{4.} Normalnya sama, $(1,1,1)$, tetapi $\frac52 \neq 1$:
jadi kedua bidangnya sejajar secara tegas. Garisnya: arahnya memberi
$(1,1,0) \cdot (1,1,1) = 2 \neq 0$, jadi tidak sejajar dengan bidangnya,
sehingga garis itu menembusnya di satu titik.

\textbf{5.} \omterm{prop:g10:coordgeom:midpoint}{Titik tengahnya}: $\left(1, \frac12, 0\right)$,
$\left(\frac12, 1, 0\right)$, $\left(0, 1, \frac12\right)$,
$\left(0, \frac12, 1\right)$, $\left(\frac12, 0, 1\right)$,
$\left(1, 0, \frac12\right)$: masing-masing berjumlah \omterm{def:g10:coordgeom:system}{koordinat}
$\frac32$.

\textbf{6.} Normal $P$ adalah $(1,1,1)$, yaitu arah
$(AG)$: jadi tegak lurus. Pusatnya,
$\left(\frac12,\frac12,\frac12\right)$, berjumlah \omterm{def:g10:coordgeom:system}{koordinat}
$\frac32$, jadi terletak pada $P$.

\textbf{7.} \omterm{prop:g10:coordgeom:midpoint}{Titik tengah} yang berurutan berselisih \omterm{def:g11:vect:vector}{vektor} seperti
$\left(-\frac12, \frac12, 0\right)$, yang panjangnya
$\frac{\sqrt2}{2}$ setiap kalinya --- jadi enam sisi yang sama panjang.

\textbf{8.} Di $\left(\frac12, 1, 0\right)$: \omterm{def:g11:vect:vector}{vektor} rusuknya
$\left(\frac12, -\frac12, 0\right)$ dan
$\left(-\frac12, 0, \frac12\right)$, dengan \omterm{def:g12:space:dot}{hasil kali skalar} $-\frac14$
dan \omterm{def:g11:scal:dot}{norma} $\frac{\sqrt2}{2}$, jadi $\cos = -\frac12$ dan sudutnya
$120^\circ$. Sisinya sama, sudut $120^\circ$-nya sama: jadi segi enam
beraturan.

\textbf{9.} Kelilingnya $6 \times \frac{\sqrt2}{2} = 3\sqrt2
\approx 4.24$. Luasnya $6 \times \frac{\sqrt3}{4} \times
\frac12 = \frac{3\sqrt3}{4} \approx 1.30$: lebih besar daripada satu sisinya
(yang berluas $1$) --- sebuah \omterm{def:g10:numbers:interunion}{irisan} yang lebih besar daripada sisi kotaknya; hanya
persegi panjang diagonalnya ($\sqrt2 \approx 1.41$) yang mengalahkannya.

\textbf{10.} Garis $(t,t,t)$ menemui $P$ di $t = \frac12$, yaitu pusatnya
--- dan itulah \omterm{prop:g10:coordgeom:midpoint}{titik tengah} $[AG]$ (sebab $A$ di $t = 0$ dan $G$ di
$t = 1$).

\textbf{11.} Untuk $c = \frac12$: bidangnya memotong ketiga rusuk
di sudut $A$, yaitu di $\left(\frac12,0,0\right)$,
$\left(0,\frac12,0\right)$, $\left(0,0,\frac12\right)$: jadi
segitiga sama sisi bersisi $\frac{\sqrt2}{2}$. Ketika $c$ tumbuh,
segitiganya membesar, sudutnya terpotong menjadi segi enam
(yang beraturan tepat di $c = \frac32$), lalu gambarnya menyusut
secara setangkup kembali menjadi segitiga di sudut $G$: itulah
perubahan bentuk mirip potongan kristal.

\textbf{12.} Tiap sisi \omterm{def:g10:numbers:interunion}{irisannya} adalah perpotongan
bidang pengirisnya dengan satu \emph{sisi} kubusnya, dan sebuah
bidang menemui satu sisi (segi banyak yang rata) paling banyak pada satu ruas:
jadi paling banyak $6$ sisi. Tujuh mustahil; tiga sampai enam semuanya
terjadi (pertanyaan 8 dan 11 memperlihatkan dua di antaranya).

\textbf{13.} Alas $ABD$: segitiga siku-siku berluas $\frac12$;
tingginya $AE = 1$, jadi volumenya
$\frac13 \times \frac12 \times 1 = \frac16$.

\textbf{14.} Tiap pasangan di antara $B, D, E, G$ berselisih tepat pada
dua \omterm{def:g10:coordgeom:system}{koordinat} sebesar $1$: jadi keenam jaraknya sama, $\sqrt2$ ---
beraturan. Kubusnya terbelah menjadi $BDEG$ ditambah empat bidang empat
sudut yang sebangun dengan $ABDE$, sehingga volumenya
$1 - 4 \times \frac16 = \frac13$.

\textbf{15.} Bidang $(BDE)$: $x + y + z = 1$; pusatnya
$\left(\frac12,\frac12,\frac12\right)$, jadi jaraknya
$\frac{\abs{\frac32 - 1}}{\sqrt3} = \frac{1}{2\sqrt3} =
\frac{\sqrt3}{6}$.

\textbf{16.} $\vect{AG} = (1,1,1)$,
$\vect{BH} = (-1,1,1)$, jadi $\cos\theta = \frac{-1 + 1 + 1}{3} =
\frac13$ dan $\theta \approx 70.5^\circ$, dengan pelurus
$\approx 109.5^\circ$. Empat pasangan elektron saling menolak
lalu menyebar sejauh mungkin di sekeliling karbonnya, dan
optimumnya adalah sudut bidang empat beraturan, yang
arah pusat-ke-titik-sudutnya bertemu tepat pada sudut ini ---
jadi intan adalah pertanyaan 14 yang mengkristal.

\textbf{17.} Titik $(1, t, t)$ pada $P$: $1 + 2t = \frac32$, jadi
$t = \frac14$ dan titiknya $\left(1, \frac14, \frac14\right)$.

\textbf{18.} Garis $(AB)$: titiknya $(t, 0, 0)$; garis $(EG)$: titiknya
$(s, s, 1)$. Agar bertemu diperlukan $0 = 1$ pada \omterm{def:g10:coordgeom:system}{koordinat}
ketiganya: mustahil. Agar sejajar diperlukan $(1,0,0)$ dan
$(1,1,0)$ \omterm{def:g12:space:collinear}{segaris}: tidak. Tidak bertemu dan tidak sejajar: jadi bersilangan
--- dua lorong pada lantai yang berbeda.

\textbf{19.} Titik $B + t(1,1,1) = (1 + t, t, t)$ pada $P$:
$1 + 3t = \frac32$, jadi $t = \frac16$ dan proyeksinya
$\left(\frac76, \frac16, \frac16\right)$, yang jumlah \omterm{def:g10:coordgeom:system}{koordinatnya}
$\frac32$: jadi pada $P$, seperti yang dituntut.

\textbf{20.} Garis parametrik menembus (pertanyaan 1, 10, 17);
\omterm{def:g12:space:normal}{vektor normal} mengukur sudut dan jarak (pertanyaan 2, 8,
15, 16); \omterm{def:g10:algebra:equation}{persamaan} bidang memahat \omterm{def:g10:numbers:interunion}{irisan} (perubahan bentuk dari
segitiga ke segi enam). Dan tempat bermainnya membalas kunjungan itu: sebuah
segi enam beraturan di dalam kotak berisi persegi, sebuah bidang empat beraturan di
sudutnya, sudut intan, dan dua garis yang saling mengabaikan
--- geometri yang hanya dapat ditawarkan oleh ruang.

\end{solution}
